\section{Objetivos}
\label{sec:objectives}

La formación de coaliciones solo resulta útil si conduce a un transporte
ejecutable. Por esta razón, los objetivos abarcan la decisión estratégica, las
restricciones físicas, los fallos y la comunicación.

\subsection{Objetivo general}

Desarrollar y evaluar una arquitectura híbrida e interpretable que combine
decisiones estratégicas locales y comunicación vecinal con cierres enteros y
guardias físicas para formar coaliciones de robots móviles autónomos y transportar
cargas heterogéneas. Medir su factibilidad, calidad de la solución, respuesta ante
perturbaciones y coste de información frente a referencias globales.

\subsection{Objetivos específicos}

\begin{enumerate}[
  label=\textbf{OE\arabic*.},
  leftmargin=3.4em,
  topsep=0.4em,
  itemsep=0.55em,
  parsep=0pt
]

  \item Caracterizar la formación de coaliciones en el caso homogéneo uno a uno,
  con cuotas variables y con servicio operacional heterogéneo. Determinar las
  condiciones de realizabilidad, cierre entero y equilibrio, y medir la pérdida
  de calidad frente a referencias globales bajo abundancia y escasez.

  \item Construir un certificado de factibilidad mecánica para la modalidad
  Cargo planar soportada a partir de la geometría de contacto, la matriz de
  agarre, los límites de actuación y el residual normalizado de \emph{wrench}.
  Vincular el certificado con el acoplamiento de robots diferenciales y con el
  servocontrol de pose dentro del dominio donde se acredita estabilidad local.

  \item Medir la respuesta de las coaliciones ante obstáculos, un fallo simple
  y conflictos de circulación. Registrar colisiones, progreso, bloqueos,
  restauración del certificado y tiempo de recuperación. La exclusión lógica
  del tráfico discreto se evaluará por separado de la seguridad física continua.

  \item Determinar cómo afectan el tamaño del sistema, el alcance limitado de la
  comunicación, el retardo y la pérdida de paquetes a la observabilidad, la
  calidad estratégica, el coste computacional y el volumen de mensajes. Comparar
  la coordinación vecinal por eventos y la retransmisión periódica con información
  perfecta y con un oráculo exhaustivo, usado solo cuando certifique su solución.

  \item Integrar reclutamiento, cierre de coalición, acoplamiento, transporte,
  seguridad y reemplazo tras un fallo en una misión Cargo. Usar ablaciones para
  identificar los componentes necesarios y acotar la transferencia de resultados
  al piloto industrial geométrico y cinemático.

\end{enumerate}
